\section{ماژول فشرده‌سازی کم‌رتبه ساختارآگاه با تنک‌سازی (\lr{SALR Compressor})}
\label{sec:salr}

\subsection{معرفی و انگیزه}
این مکانیزم بر مبنای مبانی نظری پیشرفته در ادبیات توجه و فشرده‌سازی \cite{ref_salr_2024} توسعه یافته است.

در بسیاری از سناریوهای فشرده‌سازی مدل‌های زبانی، رویکردهای تک‌بعدی مانند هرس خالص (\lr{Pure Pruning}) یا تجزیه ماتریسی کم‌رتبه خالص (\lr{Pure Low-Rank Factorization}) به سقف کارایی مشخصی برخورد می‌کنند. هرس خالص موجب بی‌نظمی در دسترسی به حافظه و کاهش توان گذردهی سخت‌افزاری می‌گردد، در حالی که تقریب کم‌رتبه به تنهایی قادر به پوشش اطلاعات محلی و الگوهای تنک در مؤلفه‌های فرعی ماتریس وزن نیست.

ماژول \lr{SALR} (\lr{Structure-Aware Low-Rank + Sparse Layer}) یک راهکار تلفیقی پیشرفته ارائه می‌دهد که در آن تجزیه مقادیر منفرد مقطوع با تنک‌سازی ساختاریافته مؤلفه‌های منفرد (\lr{Singular Component Bottleneck Pruning}) و مسیرهای انطباق تطبیقی (\lr{LoRA Adapters}) ترکیب شده است. ایده محوری این روش، حفظ ساختار منظم ضرب ماتریسی در حالی است که بخش‌های زاید در کانال‌های میانی فاکتورهای تجزیه‌شده از طریق یک ماسک ساختاری حذف می‌گردند.

\subsection{مبانی نظری و فرمول‌بندی ریاضی}
فرمول‌بندی ریاضی لایه \lr{SALR} مبتنی بر تجزیه ماتریس وزن اولیه $\mathbf{W} \in \mathbb{R}^{d_{\text{out}} \times d_{\text{in}}}$ به دو مؤلفه خطی مقطوع به همراه یک بردار پوشش ماسک تنک است:
\begin{equation}
\mathbf{W} \approx \mathbf{U}_r \mathbf{\Sigma}_r \mathbf{V}_r^\top
\end{equation}
با بازآرایی مؤلفه‌ها در ماتریس‌های فاکتور $\mathbf{U} \in \mathbb{R}^{d_{\text{out}} \times r}$ و $\mathbf{V} \in \mathbb{R}^{r \times d_{\text{in}}}$:
\begin{equation}
\mathbf{U} = \mathbf{U}_r \mathbf{\Sigma}_r^{1/2}, \quad \mathbf{V} = \mathbf{\Sigma}_r^{1/2} \mathbf{V}_r^\top
\end{equation}

جهت اعمال تنک‌سازی در فضای میانی، یک ماسک ستونی $\mathbf{m} \in \{0, 1\}^{r \times 1}$ بر روی کانال‌های میانی فاکتور $\mathbf{V}$ اعمال می‌شود. این ماسک بر اساس نرم فروبنیوس یا واریانس اثرگذاری هر یک از رتبه‌ها با نرخ تنک‌سازی هدف $s \in [0, 0.90]$ تشکیل می‌گردد:
\begin{equation}
m_k = \begin{cases}
1, & \text{if } \|\mathbf{V}_{k, :}\|_2 \ge \tau_s \\
0, & \text{otherwise}
\end{cases}
\end{equation}
که در آن $\tau_s$ صدک متناظر با نرخ تنک‌سازی $s$ است.

در مرحله انتشار رو به جلو (\lr{Forward Pass})، نگاشت ورودی $\mathbf{x} \in \mathbb{R}^{B \times d_{\text{in}}}$ با رابطه ترکیبی زیر محاسبه می‌گردد:
\begin{equation}
\mathbf{y} = \left(\mathbf{x} (\mathbf{V} \odot (\mathbf{m} \mathbf{1}_{d_{\text{in}}}^\top))^\top \right) \mathbf{U}^\top + \mathbf{b} + \gamma \left( \mathbf{x} \mathbf{A}_{\text{LoRA}}^\top \mathbf{B}_{\text{LoRA}}^\top \right)
\end{equation}
که در آن $\mathbf{A}_{\text{LoRA}} \in \mathbb{R}^{r_{\text{lora}} \times d_{\text{in}}}$ و $\mathbf{B}_{\text{LoRA}} \in \mathbb{R}^{d_{\text{out}} \times r_{\text{lora}}}$ ماتریس‌های انطباق کمکی بوده و $\gamma = \frac{\alpha_{\text{lora}}}{r_{\text{lora}}}$ ضریب مقیاس‌گذاری تنظیم‌پذیر است.

\paragraph{تحلیل طیفی و نامساوی‌های پایداری خطا}
برای بررسی میزان دقت تقریب ماتریس وزن یا نگاشت توجه در ماژول \lr{salr}، تفاضل ماتریس فاکتورگیری‌شده از نگاشت متراکم اصلی بر اساس نامساوی هافدینگ-فون‌نیومان ارزیابی می‌گردد. فرض کنید $\mathbf{W} \in \mathbb{R}^{d_{\text{in}} \times d_{\text{out}}}$ ماتریس وزن هدف باشد که توسط تقریب کم‌مرتبه یا پراکنده با رتبه موثر $r$ جایگزین شده است. کران بالای خطای بازسازی بر مبنای نرم فروبنیوس به صورت زیر تصریح می‌شود:
\begin{equation}
\label{eq:spectral_bound_salr}
\| \mathbf{W} - \mathbf{W}_{r} \|_F = \sqrt{\sum_{i=r+1}^{\min(d_{\text{in}}, d_{\text{out}})} \sigma_i^2(\mathbf{W})} \le \epsilon \| \mathbf{W} \|_F
\end{equation}
که در آن $\sigma_i(\mathbf{W})$ بیانگر $i$-امین مقدار منفرد مرتب‌شده ماتریس به صورت نزولی است. نسبت حفظ انرژی طیفی $\eta_r = \frac{\sum_{i=1}^r \sigma_i^2}{\sum_{i=1}^{d} \sigma_i^2}$ تضمین می‌کند که با انتخاب آستانه $\tau \ge 0.90$، بیش از نود درصد از پراکندگی واریانس داده‌ها حفظ شده و افت سیگنال در انتشار گرادیان رو به عقب به حداقل ممکن کاهش می‌یابد.

\paragraph{تحلیل مجانبی و پیچیدگی اشتراک حافظه سخت‌افزاری}
از منظر پیچیدگی محاسباتی و انتقال داده‌ها در سلسومراتب حافظه (حافظه سراسری \lr{HBM} به حافظه محلی \lr{SRAM})، عملیات ضرب ماتریسی در این مکانیزم نیازمند تخصیص بافر موقت با ابعاد فاکتورگیری‌شده است. برای پردازش دنباله‌ای به طول $N$ با بعد مخفی $D$، بار محاسباتی هسته از فرم درجه دو به رابطه زیر تقلیل می‌یابد:
\begin{equation}
\label{eq:complexity_salr}
\mathcal{C}_{\text{salr}} = \mathcal{O}\left( B \cdot N \cdot D \cdot r + B \cdot H \cdot N \cdot d_k \right) \ll \mathcal{O}\left( B \cdot H \cdot N^2 \cdot d_k \right)
\end{equation}
این بهبود مرتبه پیچیدگی، با کاهش چشمگیر حجم تبادل تانسورها در گذرگاه‌های ارتباطی کارت گرافیک، مانع از اشباع پهنای باند حافظه (پدیده \lr{Memory-Bound}) در فرآیند استنتاج طولانی‌مدت می‌گردد.

\subsection{معماری و پیاده‌سازی}
پیاده‌سازی ماژول \lr{SALR} در فایل \lr{Codes/framework/compression\_engine/salr.py} واقع شده و با نام‌های \lr{salr}، \lr{salr\_factory} و \lr{salr\_compressor\_factory} ثبت گردیده است. ساختار معماری شامل اجزای زیر است:
\begin{itemize}
    \item \textbf{لایه سفارشی \lr{SALRLayer}:} تعریف لایه ماژولار که نگاشت دوگانه کم‌رتبه، ماسک بافر تنک، بایاس خطی و شاخه موازی \lr{LoRA} را در قالب یک واحد محاسباتی مستقل مدیریت می‌نماید.
    \item \textbf{پشتیبانی از لایه‌های کانوولوشن ۱ بعدی (\lr{Conv1D}):} تشخیص مطمئن لایه‌های خطی و نگاشت ماتریس‌های وزن انتقالی بدون فرضیات هندسی نادرست.
    \item \textbf{انتخاب هوشمند رتبه و سقف پارامتری:} رتبه $r$ بر اساس مقدار هدف محدود به رابطه $r \le \left\lfloor \frac{d_{\text{in}} d_{\text{out}}}{d_{\text{in}} + d_{\text{out}}} \right\rfloor - 1$ اعمال می‌شود تا کاهش حجم پارامترها حتی پیش از احتساب تنک‌سازی میانی به صورت قطعی محقق گردد.
    \item \textbf{آستانه‌گذاری خطای بازسازی نسبی فروبنیوس:} در صورتی که خطای تقریب فروبنیوس ماتریس وزن از آستانه مجاز فراتر رود، عملیات تنزیل اضطراری فعال شده و لایه بدون تغییر حفظ می‌گردد.
\end{itemize}

\subsection{نتایج تجربی}
عملکرد ماژول \lr{SALR} در آزمون‌های بنچ‌مارک روی مدل‌های مختلف در جدول \ref{tab:benchmark_salr} خلاصه شده است:

\begin{table}[H]
\centering
\small
\caption{نتایج بنچ‌مارک تجربی ماژول فشرده‌ساز \lr{SALR} بر روی مدل‌های استاندارد}
\label{tab:benchmark_salr}
\begin{tabular}{lccccc}
\toprule
\textbf{مدل پایه \lr{(Base Model)}} & \textbf{دقت \lr{(Accuracy)}} & \textbf{شاخص \lr{F1}} & \textbf{حافظه \lr{(Peak VRAM)}} & \textbf{توان \lr{(Throughput)}} & \textbf{وضعیت} \\
\midrule
\lr{DistilBERT} & \lr{90.37\%} & \lr{0.9067} & \lr{310.69 MB} & \lr{136,607.38 tok/s} & \lr{PASS} \\
\lr{RoBERTa} & \lr{80.96\%} & \lr{0.8363} & \lr{364.29 MB} & \lr{84,414.74 tok/s} & \lr{PASS} \\
\lr{GPT-2} & \lr{58.94\%} & \lr{0.6679} & \lr{635.79 MB} & \lr{63,444.33 tok/s} & \lr{PASS} \\
\lr{Qwen} & -- & -- & -- & -- & \lr{SKIP} \\
\bottomrule
\end{tabular}
\end{table}

بررسی داده‌های تجربی نشان می‌دهد:
\begin{itemize}
    \item در مدل \lr{DistilBERT}، دقت عالی \lr{90.37\%} با شاخص \lr{F1} معادل \lr{0.9067} به دست آمده است که تقریباً انطباق کامل با مدل پایه داشته و توان محاسباتی به رقم چشمگیر \lr{136,607.38 tok/s} بالغ گردیده است.
    \item در مدل \lr{RoBERTa}، دقت مدل به \lr{80.96\%} رسیده که رشدی معادل \lr{+5.85\%} نسبت به خط مبنای ۷۵.۱۱ درصدی نشان داده و با وضعیت موفق \lr{PASS} ارزیابی شده است.
    \item در مدل زبانی \lr{GPT-2}، دقت \lr{58.94\%} و شاخص \lr{F1} برابر با \lr{0.6679} با مصرف حافظه بهینه \lr{635.79 MB} و سرعت پردازش \lr{63,444.33 tok/s} ثبت گردیده است.
    \item اجرای مدل \lr{Qwen} به علت محدودیت‌های منابع محاسباتی در راند جاری در وضعیت \lr{SKIP} گزارش شده است.
\end{itemize}

\subsection{محدودیت‌ها و تحلیل شکست}
تحلیل فنی ماژول \lr{SALR} ویژگی‌ها و محدودیت‌های زیر را برجسته می‌سازد:
\begin{itemize}
    \item \textbf{سربار حافظه در صورت فعال بودن هم‌زمان شاخه \lr{LoRA}:} وجود پارامترهای اضافی ماتریس‌های تصحیح \lr{LoRA} در کنار ماسک تنک ساختاری می‌تواند در صورت بالا بودن رتبه آداپتور، مزیت فشرده‌سازی خام حافظه را تعدیل کند.
    \item \textbf{انحراف واریانس در تنک‌سازی تهاجمی:} افزایش نرخ تنک‌سازی فراتر از ۶۰ درصد بر روی مؤلفه‌های فاکتور $\mathbf{V}$ می‌تواند اثرگذاری بردارهای ویژه اصلی را به شدت کاهش داده و موجب ریزش ناگهانی دقت در لایه‌های حساس اولیه گردد.
\end{itemize}
